\section{Gradient Estimators}
\label{sec:methods}
Consider one scalar parameter $\lambda$, either a mixture logit or a coordinate
of a component parameter. We write $\widehat g_\lambda$ for an estimate of
$\partial_\lambda\mathcal J$ in Eq.~\eqref{eq:objective}.
The derivative $\partial_\lambda f_\theta(z)$ holds the sample point fixed.
If the point moves as $\lambda$ changes, its contribution to the objective
also depends on the slope $\nabla_z f_\theta(z)$.

\paragraph*{Score-function.}
Draw $Z\sim q_\theta$ and hold it fixed. Differentiating the density in the
expectation gives $\partial_\lambda q_\theta(z)
=q_\theta(z)\,\partial_\lambda\log q_\theta(z)$. The score therefore measures
the relative change in density at the sampled point~\citep{ranganath2014bbvi}.
With a constant baseline $b$, the estimate is
\begin{equation}
 \widehat g_\lambda
 = \partial_\lambda f_\theta(Z)
   + \bigl(f_\theta(Z)-b\bigr)\,\partial_\lambda\log q_\theta(Z).
 \label{eq:score}
\end{equation}
The score has mean zero, so $b$ can center the sampled objective values
while preserving the expected gradient.

\paragraph*{Stratified.}
Use the component sum in Eq.~\eqref{eq:objective} and draw one independent
$Z_k=T_k(\varepsilon_k;\phi_k)$ per component~\citep{morningstar2021mixtures}. Write
$v_k(e)=\partial_\lambda T_k(e;\phi_k)$ for the sample derivative at fixed
noise $e$. Differentiating the weighted sample objective
$\sum_k\pi_k f_\theta(Z_k)$ gives
\begin{equation}
 \widehat g_\lambda
 = \sum_{k=1}^K \left[
     f_\theta(Z_k)\,\partial_\lambda\pi_k
     + \pi_k\bigl\{\partial_\lambda f_\theta(Z_k)
          + \nabla_z f_\theta(Z_k)^\top v_k(\varepsilon_k)\bigr\}\right].
 \label{eq:stratified}
\end{equation}
The weight term for logit $\alpha_j$ is positive when $f_\theta(Z_j)$ exceeds
the weighted sample average. The path term uses the objective slope at
each component's own draw.

\paragraph*{Implicit reparameterization.}
Draw $Z\sim q_\theta$. In one dimension the mixture cumulative distribution
function (CDF) $F_\theta$ assigns it a rank $U=F_\theta(Z)$, uniform on
$[0,1]$. As $\lambda$ varies, follow the point with that same
rank~\citep{figurnov2018implicit}. Differentiating $F_\theta(Z)=U$ and
applying the chain rule gives the sample velocity $v(Z)$ and the estimate
\begin{equation}
 v(Z)=-\frac{\partial_\lambda F_\theta(Z)}{q_\theta(Z)},\qquad
 \widehat g_\lambda
 = \partial_\lambda f_\theta(Z)+\nabla_z f_\theta(Z)^\top v(Z).
 \label{eq:implicit}
\end{equation}
If $F_j$ denotes the CDF of component $j$, then
$\partial_{\alpha_j}F_\theta=\pi_j(F_j-F_\theta)$ gives the numerator
for a logit. In higher dimensions, hold fixed the conditional ranks
$S_{\theta,i}(z)=F_\theta(z_i\mid z_1,\ldots,z_{i-1})$ for $i=1,\ldots,d$.
Their Jacobian $D_z S_\theta$ is lower triangular. Solving
$(D_z S_\theta)v=-\partial_\lambda S_\theta$ at $Z$ gives the velocity in
Eq.~\eqref{eq:implicit} wherever the conditional densities are positive.

\paragraph*{Post-stratified.}
Draw one $Z\sim q_\theta$. For a given label $C=k$, a component parameter
enters through the sampling map $T_k$ and a logit through the probability
$\pi_k$ of selecting $k$. Post-stratification averages these contributions
over the possible labels of $Z$ with probabilities $r_k(Z)$~\citep{burroni2026mixture}.
If each component map is invertible, the noise $e_k=T_k^{-1}(Z;\phi_k)$
reconstructs the same point in component $k$. Using these noises gives
\begin{equation}
 \widehat g_\lambda
 = \partial_\lambda f_\theta(Z)
   + \sum_{k=1}^K r_k(Z)\left[
       \nabla_z f_\theta(Z)^\top v_k(e_k)
       + f_\theta(Z)\,\partial_\lambda\log\pi_k\right].
 \label{eq:poststratified}
\end{equation}
For a univariate Gaussian component with mean $\mu_k$ and scale $\sigma_k$,
$e_k=(Z-\mu_k)/\sigma_k$. Holding $e_k$ fixed gives a unit sample derivative
for $\mu_k$, so its sampling contribution is $r_k(Z)\,\partial_z f_\theta(Z)$.
